\begin{ex}
Rút gọn phân số $\frac{15}{25}$ về tối giản.
\shortans{3/5}
\loigiai{
$\frac{3}{5}$.
}
\end{ex}

\begin{ex}
Rút gọn phân số $\frac{18}{24}$ về tối giản.
\shortans{3/4}
\loigiai{
$\frac{3}{4}$.
}
\end{ex}

\begin{ex}
Rút gọn phân số $\frac{20}{30}$ về tối giản.
\shortans{2/3}
\loigiai{
$\frac{2}{3}$.
}
\end{ex}

\begin{ex}
Rút gọn phân số $\frac{36}{54}$ về tối giản.
\shortans{2/3}
\loigiai{
$\frac{2}{3}$.
}
\end{ex}

\begin{ex}
Quy đồng mẫu số $\frac{2}{5}$ và $\frac{3}{10}$ với mẫu chung 10, phân số $\frac{2}{5}$ trở thành phân số nào?
\shortans{4/10}
\loigiai{
$\frac{4}{10}$.
}
\end{ex}

\begin{ex}
Mẫu số chung nhỏ nhất của $\frac{1}{6}$ và $\frac{3}{8}$ là bao nhiêu?
\shortans{24}
\loigiai{
24.
}
\end{ex}

\begin{ex}
So sánh $\frac{4}{7}$ và $\frac{5}{7}$: phân số nào lớn hơn?
\shortans{5/7}
\loigiai{
$\frac{5}{7}$.
}
\end{ex}

\begin{ex}
So sánh $\frac{3}{4}$ và $\frac{3}{5}$: phân số nào lớn hơn?
\shortans{3/4}
\loigiai{
$\frac{3}{4}$.
}
\end{ex}

\begin{ex}
Trong các phân số $\frac{5}{6}, \frac{6}{6}, \frac{7}{6}$, phân số nào bằng 1?
\shortans{6/6}
\loigiai{
$\frac{6}{6} = 1$.
}
\end{ex}

\begin{ex}
Phân số nào bé hơn 1 trong các phân số trên?
\shortans{5/6}
\loigiai{
$\frac{5}{6} < 1$.
}
\end{ex}

\begin{ex}
Phân số nào lớn hơn 1 trong các phân số trên?
\shortans{7/6}
\loigiai{
$\frac{7}{6} > 1$.
}
\end{ex}

\begin{ex}
Tìm $x$, biết: $\frac{x}{15} = \frac{2}{3}$.
\shortans{10}
\loigiai{
$x = 2 \times 5 = 10$.
}
\end{ex}

\begin{ex}
Tìm $x$, biết: $\frac{12}{x} = \frac{3}{4}$.
\shortans{16}
\loigiai{
$x = 4 \times 4 = 16$.
}
\end{ex}

\begin{ex}
Viết số 4 dưới dạng phân số có mẫu số bằng 3.
\shortans{12/3}
\loigiai{
$4 = \frac{12}{3}$.
}
\end{ex}

\begin{ex}
Viết số 1 dưới dạng phân số có tử số bằng 9.
\shortans{9/9}
\loigiai{
$1 = \frac{9}{9}$.
}
\end{ex}

\begin{ex}
Có bao nhiêu phân số tối giản trong các phân số: $\frac{2}{5}, \frac{4}{8}, \frac{3}{7}, \frac{6}{9}$?
\shortans{2}
\loigiai{
$\frac{2}{5}$ và $\frac{3}{7}$ (có 2 phân số).
}
\end{ex}

\begin{ex}
So sánh $\frac{2}{3}$ và $\frac{4}{6}$: hai phân số có bằng nhau không? (Bằng nhau ghi 1, khác nhau ghi 0).
\shortans{1}
\loigiai{
Bằng nhau (1).
}
\end{ex}

\begin{ex}
Phân số chỉ 45 phút so với 1 giờ rút gọn về tối giản là bao nhiêu?
\shortans{3/4}
\loigiai{
$\frac{45}{60} = \frac{3}{4}$.
}
\end{ex}

\begin{ex}
Phân số chỉ 20 cm so với 1 m rút gọn về tối giản là bao nhiêu?
\shortans{1/5}
\loigiai{
$\frac{20}{100} = \frac{1}{5}$.
}
\end{ex}

\begin{ex}
Phân số chỉ 250 g so với 1 kg rút gọn về tối giản là bao nhiêu?
\shortans{1/4}
\loigiai{
$\frac{250}{1000} = \frac{1}{4}$.
}
\end{ex}

\begin{ex}
Tử số của phân số $\frac{8}{15}$ là bao nhiêu?
\shortans{8}
\loigiai{
Tử số là 8.
}
\end{ex}

\begin{ex}
Mẫu số của phân số $\frac{8}{15}$ là bao nhiêu?
\shortans{15}
\loigiai{
Mẫu số là 15.
}
\end{ex}

\begin{ex}
Phân số $\frac{0}{12}$ có giá trị bằng bao nhiêu?
\shortans{0}
\loigiai{
Bằng 0.
}
\end{ex}

\begin{ex}
Trong hai phân số $\frac{7}{9}$ và $\frac{8}{9}$, phân số lớn hơn là phân số nào?
\shortans{8/9}
\loigiai{
$\frac{8}{9}$.
}
\end{ex}

\begin{ex}
Tìm $y$: $\frac{y}{20} = \frac{1}{4}$.
\shortans{5}
\loigiai{
$y = 20 : 4 = 5$.
}
\end{ex}
